%! AP Statistics — Unit 1, Lesson 1.6 — Student Packet Cover Sheet
\documentclass[10pt]{article}
\usepackage{apstats-article}
\usepackage{apstats-boxes}

%=============================================================================
\begin{document}

% ── Full-bleed navy banner ────────────────────────────────────────────────────
\begin{tikzpicture}[remember picture, overlay]
  \fill[navy] (current page.north west)
    rectangle ([yshift=-0.9in]current page.north east);
\end{tikzpicture}
\vspace{-0.8in}

\begin{center}
  {\color{white}\textbf{\LARGE AP Statistics}}\\[4pt]
  {\color{skymid}\large\textbf{Unit 1: Exploring One-Variable Data}}\\[6pt]
  {\color{white}\textbf{\large Lesson 1.6 \quad Describing the Distribution of a Quantitative Variable}}
\end{center}

\vspace{0.22in}

\namedateperiod

\vspace{0.18in}

% ── LEARNING TARGETS ─────────────────────────────────────────────────────────
\begin{learningtargetbox}
\small
\begin{enumerate}[leftmargin=1.5em, itemsep=5pt]
  \item \ldots{} identify and name the shape of a distribution (symmetric,
        skewed right, skewed left, unimodal, bimodal, uniform) from a histogram,
        dotplot, or stemplot.
  \item \ldots{} write a complete description of a quantitative distribution
        using the \textbf{SOCS} framework: Shape, Outliers, Center, and Spread.
  \item \ldots{} include context (variable name and population) in every sentence
        of a SOCS description.
  \item \ldots{} explain why outliers are important and how they affect a
        reader's interpretation of center and spread.
\end{enumerate}
\end{learningtargetbox}

\vspace{0.14in}

% ── TABLE OF CONTENTS ────────────────────────────────────────────────────────
\begin{tocbox}
\small
\renewcommand{\arraystretch}{1.5}
\begin{tabularx}{\linewidth}{@{} c l X r @{}}
\rowcolor{sky}
\textbf{\#} & \textbf{Component} & \textbf{Description} & \textbf{Score} \\
\midrule
1 & Warm-Up        & Shape vocabulary + reading a distribution display          & \blank{1.2cm} \\
2 & Guided Notes   & SOCS framework, shape vocabulary, worked examples           & \blank{1.2cm} \\
3 & Group Activity & Write and critique complete SOCS descriptions               & \blank{1.2cm} \\
4 & Exit Ticket    & Write a SOCS description from a described distribution      & \blank{1.2cm} \\
5 & Homework       & SOCS descriptions, shape identification, AP free-response   & \blank{1.2cm} \\
\midrule
  & \multicolumn{2}{r}{\textbf{Total}} & \blank{1.2cm} \\
\end{tabularx}
\end{tocbox}

\vspace{0.14in}

% ── KEEP IN MIND ─────────────────────────────────────────────────────────────
\begin{remindbox}
\small
SOCS is the AP standard for describing any quantitative distribution.
A response missing even one component --- or lacking context --- cannot earn
full credit on the AP exam.

\vspace{6pt}
\begin{tabularx}{\linewidth}{@{} >{\centering\arraybackslash\columncolor{goldbg}}p{0.06\linewidth}
                                    X @{}}
\textbf{\large!} &
\textbf{Every SOCS sentence needs context.}\enspace
Write ``The distribution of SAT scores for these 200 students\ldots'' not
just ``The distribution is\ldots'' AP graders look for context in every
sentence, not just the first one. \\[6pt]
\textbf{\large!} &
\textbf{Skewed right means the tail points right.}\enspace
The peak (bulk of data) is on the left; a few unusually large values pull
the tail toward higher values. Do \emph{not} confuse the peak location with
the direction of skew. \\[6pt]
\textbf{\large!} &
\textbf{Outliers affect center and spread.}\enspace
Always mention potential outliers and note that they pull the mean away from
the median. Investigate the cause; do not automatically remove them. \\[6pt]
\textbf{\large!} &
\textbf{Center and spread are informal here.}\enspace
In Lesson 1.6 we estimate center and spread by eye. Lesson 1.7 gives us
exact numerical measures (mean, median, IQR, standard deviation). \\
\end{tabularx}
\end{remindbox}

\end{document}
